% !PRINTED 188
\weberpage{188}

Wenn wir nach den \S.\,16, (4) die Function
\[
F(x_1 + t\xi_1, x_2 + t\xi_2 \ldots x_m + t\xi_m)
\]
nach Potenzen von $t$ ordnen, so erhalten wir für den Coëfficienten von $t^\nu$
\weq{4}{F_\nu(x,\xi) = \frac{1}{\Pi(\nu)}\sum^{\alpha}\frac{\Pi(\nu)}{\Pi(\alpha_1) \ldots \Pi(\alpha_m)}\xi_1^{\alpha_1}\xi_2^{\alpha_2}\ldots\xi_m^{\alpha_m}D_{\alpha_1,\alpha_2\ldots\alpha_m}(F),}
\[
\alpha_1 + \alpha_2 + \cdots + \alpha_m = \nu.
\]
Wenn wir auf die Variablen $\xi$ und $x$ gleichzeitig die lineare
Transformation (1), \S.\,55 anwenden, so geht $F_\nu(x,\xi)$ in $\Phi_\nu(x',\xi')$
über, was man aus $F_\nu$ erhält, wenn man gleichzeitig alle Buchstaben
$x, \xi, a$ durch $x', \xi', a'$ und $F$ durch $\Phi$ ersetzt.

Daraus folgt:

I. \gesperrt{Betrachtet man $F_\nu(x,\xi)$ als Form $\nu^{\mathrm{ten}}$ Grades
von $\xi$ und bildet eine Invariante dieser Form, deren Coëfficienten also
noch von den $x$ abhängen, so erhält man eine Covariante der Form $F$.}

Und da die Variablen $x$ und $\xi$ derselben Transformation
unterliegen, so ergiebt sich ebenso:

II. \gesperrt{Bildet man eine Covariante der Function $F_\nu(x,\xi)$ als
Function von $\xi$ betrachtet, und ersetzt darin die $\xi$ durch die $x$,
so erhält man eine Covariante der ursprünglichen Function $F(x)$.}

Die gleichen Sätze gelten auch für simultane Invarianten
und Covarianten.

Von den Functionen $F_\nu(x,\xi)$ gilt der Satz
\weq{5}{F_\nu(x,\xi) = F_{n-\nu}(\xi,x),}
der sich durch die Vergleichung entsprechender Potenzen von $t$
auf beiden Seiten der Identität
\[
F(x_1+t\xi_1\ldots x_m+t\xi_m) = t^n F(t^{-1}x_1+\xi_1\ldots t^{-1}x_m+\xi_m)
\]
ergiebt.

Die Functionen
\[
P_\nu(x,\xi) = \frac{\Pi(\nu)\;\Pi(n-\nu)}{\Pi(n)}F_\nu(x,\xi)
\]
werden die \gesperrt{Polaren der Function $F(x)$} genannt, und zwar
heisst $P_\nu(x,\xi)$ die $\nu^{\mathrm{te}}$ Polare. Auch von ihnen gilt der Satz
\[
P_\nu(x,\xi) = P_{n-\nu}(\xi,x).
\]
